% 第4章 磁有序与磁结构？不——第4章 相互作用（Interactions）
\chapter{相互作用}

本章讨论几类磁相互作用：它们使固体中的磁矩得以相互“沟通”，并可能产生长程序。

\section{磁偶极相互作用}

最先想到会起作用的相互作用是磁偶极相互作用。两个相距 $\mathbf{r}$ 的磁偶极子
$\boldsymbol{\mu}_1$ 与 $\boldsymbol{\mu}_2$ 的能量为
\begin{equation*}
E = \frac{\mu_0}{4\pi r^{3}}\left[\boldsymbol{\mu}_1\cdot\boldsymbol{\mu}_2 - \frac{3}{r^{2}}(\boldsymbol{\mu}_1\cdot\mathbf{r})(\boldsymbol{\mu}_2\cdot\mathbf{r})\right]\tag{4.1}
\end{equation*}
依赖其间距与相互排列的程度。容易估计其量级：两个 $\mu\approx1\mu_{\mathrm{B}}$ 的磁矩
相距 $r\approx1$ \AA{} 时约为 $\mu^{2}/4\pi r^{3}\sim10^{-23}$~J，按温度计约 1 K。
\fn{1}{能量与温度的便捷换算见图 3.8。}许多材料在远高于此的温度有序（有的高达
1000 K 左右），因此磁偶极相互作用太弱，不足以解释大多数磁性材料的有序。尽管如此，
对在毫开尔文量级有序的那些材料，它可能相当重要。

\section{交换相互作用}

交换相互作用位于长程磁有序现象的心脏。交换效应精微而不无神秘——处理的无非是
一根条形磁铁和一堆铁屑，却必须动用交换算符与全同粒子，这乍看令人诧异。但正如
磁学中常见的情形，这恰恰表明量子力学是许多日常现象的根源。交换相互作用不过是
静电相互作用：同号电荷靠近时要付出能量、分开时省下能量。

\subsection{交换的起源}

考虑只有两个电子的简单模型，其空间坐标分别为 $\mathbf{r}_1$ 与 $\mathbf{r}_2$。
联合态的波函数可以写成单电子态的乘积：若第一个电子处于 $\psi_a(\mathbf{r}_1)$、
第二个电子处于 $\psi_b(\mathbf{r}_2)$，则联合波函数为
$\psi_a(\mathbf{r}_1)\psi_b(\mathbf{r}_2)$。但这个乘积态不满足交换对称性：交换两个
电子得到 $\psi_a(\mathbf{r}_2)\psi_b(\mathbf{r}_1)$，它与初始态不成倍数。因此我们
只能构造在粒子交换下行为正确的（反）对称化乘积态（1.3.4 节已讨论）。

对电子，总波函数必须反对称，故波函数的自旋部分要么是反对称单态 $\chi_{\mathrm{S}}$
（S=0，配对称空间态），要么是对称三重态 $\chi_{\mathrm{T}}$（S=1，配反对称空间态）。
单态 $\Psi_{\mathrm{S}}$ 与三重态 $\Psi_{\mathrm{T}}$ 的波函数写为
\begin{equation*}
\begin{aligned}
\Psi_{\mathrm{S}} &= \frac{1}{\sqrt{2}}\left[\psi_a(\mathbf{r}_1)\psi_b(\mathbf{r}_2) + \psi_a(\mathbf{r}_2)\psi_b(\mathbf{r}_1)\right]\chi_{\mathrm{S}}\\
\Psi_{\mathrm{T}} &= \frac{1}{\sqrt{2}}\left[\psi_a(\mathbf{r}_1)\psi_b(\mathbf{r}_2) - \psi_a(\mathbf{r}_2)\psi_b(\mathbf{r}_1)\right]\chi_{\mathrm{T}},
\end{aligned}\tag{4.2}
\end{equation*}
空间与自旋部分都包括在内。两个可能态的能量为
\begin{equation*}
\begin{aligned}
E_{\mathrm{S}} &= \int\Psi_{\mathrm{S}}^{*}\hat{\mathcal{H}}\Psi_{\mathrm{S}}\,\mathrm{d}\mathbf{r}_1\mathrm{d}\mathbf{r}_2\\
E_{\mathrm{T}} &= \int\Psi_{\mathrm{T}}^{*}\hat{\mathcal{H}}\Psi_{\mathrm{T}}\,\mathrm{d}\mathbf{r}_1\mathrm{d}\mathbf{r}_2,
\end{aligned}
\end{equation*}
假定波函数的自旋部分 $\chi_{\mathrm{S}}$、$\chi_{\mathrm{T}}$ 已归一化。两能量之差为
\begin{equation*}
E_{\mathrm{S}} - E_{\mathrm{T}} = 2\int\psi_a^{*}(\mathbf{r}_1)\psi_b^{*}(\mathbf{r}_2)\hat{\mathcal{H}}\psi_a(\mathbf{r}_2)\psi_b(\mathbf{r}_1)\,\mathrm{d}\mathbf{r}_1\mathrm{d}\mathbf{r}_2.\tag{4.3}
\end{equation*}
式 1.70 表明单态与三重态之差可以用 $\mathbf{S}_1\cdot\mathbf{S}_2$ 参数化：单态
$\mathbf{S}_1\cdot\mathbf{S}_2=-\frac{3}{4}$，三重态 $\mathbf{S}_1\cdot\mathbf{S}_2=\frac{1}{4}$。
于是哈密顿量可以写成“有效哈密顿量”的形式
\begin{equation*}
\hat{\mathcal{H}} = \frac{1}{4}(E_{\mathrm{S}}+3E_{\mathrm{T}}) - (E_{\mathrm{S}}-E_{\mathrm{T}})\mathbf{S}_1\cdot\mathbf{S}_2.\tag{4.4}
\end{equation*}
这是常数项与依赖自旋的项之和。常数可并入其他常数能量项，第二项更有意思。交换常数
（或交换积分）$\mathsf{J}$ 定义为
\begin{equation*}
\mathsf{J} = \frac{E_{\mathrm{S}}-E_{\mathrm{T}}}{2} = \int\psi_a^{*}(\mathbf{r}_1)\psi_b^{*}(\mathbf{r}_2)\hat{\mathcal{H}}\psi_a(\mathbf{r}_2)\psi_b(\mathbf{r}_1)\,\mathrm{d}\mathbf{r}_1\mathrm{d}\mathbf{r}_2,\tag{4.5}
\end{equation*}
于是有效哈密顿量中依赖自旋的项可写作
\begin{equation*}
\hat{\mathcal{H}}^{\text{spin}} = -2\mathsf{J}\,\mathbf{S}_1\cdot\mathbf{S}_2.\tag{4.6}
\end{equation*}
若 $\mathsf{J}>0$，$E_{\mathrm{S}}>E_{\mathrm{T}}$，更倾向于三重态 $S=1$；若
$\mathsf{J}<0$，$E_{\mathrm{S}}<E_{\mathrm{T}}$，更倾向于单态 $S=0$。对两个电子导出
此式还算简单，\bio{Werner Heisenberg\\（1901--1976）}\mnote{注意有些书把 J 取为此处数值的两倍，此时式 4.6--4.8 变为
\begin{equation*}
\hat{\mathcal{H}}^{\mathrm{spin}} = -\mathsf{J}\,\mathbf{S}_1\cdot\mathbf{S}_2,\tag{4.9}
\end{equation*}
\begin{equation*}
\hat{\mathcal{H}} = -\frac{1}{2}\sum_{ij}\mathsf{J}_{ij}\,\mathbf{S}_i\cdot\mathbf{S}_j,\tag{4.10}
\end{equation*}
\begin{equation*}
\hat{\mathcal{H}} = -\sum_{i>j}\mathsf{J}_{ij}\,\mathbf{S}_i\cdot\mathbf{S}_j.\tag{4.11}
\end{equation*}}
但推广到多体系统绝非易事。尽管如此，量子力学早期人们就已认识到：式 4.6 这类相互
作用很可能存在于所有相邻原子之间。这引出海森伯模型（Heisenberg model）的哈密顿量：
\begin{equation*}
\hat{\mathcal{H}} = -\sum_{ij}\mathsf{J}_{ij}\,\mathbf{S}_i\cdot\mathbf{S}_j,\tag{4.7}
\end{equation*}
$\mathsf{J}_{ij}$ 是第 $i$ 与第 $j$ 个自旋之间的交换常数。因子 2 被省去，因为求和把
每对自旋计了两次。式 4.7 的另一种写法是
\begin{equation*}
\hat{\mathcal{H}} = -2\sum_{i>j}\mathsf{J}_{ij}\,\mathbf{S}_i\cdot\mathbf{S}_j,\tag{4.8}
\end{equation*}
$i>j$ 避免了“重复计数”，因子 2 于是回来了。常可取 $\mathsf{J}_{ij}$ 对最近邻自旋
等于常数 $\mathsf{J}$、其余为零。

交换积分的计算一般会相当复杂，这里只提几个一般性的特征。第一，若两个电子在同一原子上，
交换积分通常为正：这稳定三重态、保证空间态反对称，使两电子相互回避从而减小库仑
排斥——与洪特第一条规则一致。

当两个电子分处相邻原子时，情形截然不同。任何联合态都是一个以某原子为中心的态与
以另一原子为中心的态的组合。值得记住：长为 $L$ 的一维箱中粒子的能量正比于
$L^{-2}$——这是动能，说明被关进小盒子要付出巨大动能。因此电子可以通过成键节省
动能：成键让它们在两个原子而非一个原子的范围内游走（即住进“更大的盒子”）。
此时该考虑的不再是原子轨道而是分子轨道（见图 4.1）：成键轨道（空间对称）与
“反键”轨道（空间反对称），后者能量更高——反键轨道曲率更大，动能更大。这更倾向于
单态（反对称），因此交换积分很可能为负。

\subsection{直接交换}

如果相邻磁性原子上的电子经由交换相互作用发生耦合，这就称为直接交换（direct
exchange）：交换相互作用直接进行，无需任何中介。尽管这看似是交换相互作用最显而易
见的途径，物理实际却很少如此简单。

在很多情形下，直接交换并不是决定磁性的重要机制，因为相邻磁性轨道之间缺乏足够的直
接交叠。例如在稀土中，4f 电子高度局域、十分贴近原子核，其概率密度明显伸出原
子间距约十分之一以外的部分微乎其微，因此直接交换相互作用在稀土中不太可能十分有效。
即便在 Fe、Co、Ni 等过渡金属中——3d 轨道离核伸得更远——也很难解释直接交换何以导
致观测到的磁性质。这些材料是金属，这意味着传导电子的作用不容忽视，正确的描述必须
同时顾及电子的局域性与能带特征。

\begin{figure}
\centering
\includegraphics[width=0.55\textwidth]{fig4_01.pdf}
\caption{双原子分子的分子轨道。成键轨道（两原子轨道之和；就波函数空间部分而言在
交换下对称）能量低于反键轨道（两原子轨道之差；交换下反对称）。因此两个电子填入
成键态、反键态空着的单态基态更稳定。此图适用于氢分子 $\mathrm{H_2}$——其能量低于
两个孤立 H 原子（$E_0$）。注意氦的双原子形式 $\mathrm{He_2}$ 不能形成：两个 He 原子
的四个电子会同时填满成键与反键轨道，相比两个孤立 He 原子没有净能量节省。}
\end{figure}

因此在许多磁性材料中，必须考虑某种间接交换相互作用。

\subsection{离子固体中的间接交换：超交换}

一些离子固体（包括某些氧化物与氟化物）具有磁基态。例如 MnO（见图 4.2）与
$\mathrm{MnF_2}$ 都是反铁磁体——初看令人惊讶，因为两个体系中 $\mathrm{Mn^{2+}}$
离子的电子并无直接交叠。交换相互作用通常是短程的，此处起作用的较长程相互作用必在
某种意义上是“超”的。

这里起作用的交换机制称为超交换（superexchange）：可定义为非相邻磁性离子之间经
由置于其间的非磁性离子传递的间接交换相互作用。它的产生是因为反铁磁性在动能上有
好处。参照图 4.3 即可理解：图中两个过渡金属离子被一个氧离子隔开。为简单起见，设
过渡金属离子上的磁矩来自单个未成对电子（更复杂的情形可类似处理）。若系统完全
离子化，每个金属离子的 d 轨道上有一个未成对电子，而氧的最外层占据态有两个 p 电子。
图示表明：反铁磁耦合使体系能量降低——它让这些电子得以在整个结构上离域化，从而降低动能。

\begin{figure}
\centering
\includegraphics[width=0.4\textwidth]{fig4_02.pdf}
\caption{MnO 的晶体结构。$\mathrm{Mn^{2+}}$（锰）离子的最近邻对经由
$\mathrm{O^{2-}}$（氧）离子连接。}
\end{figure}

\begin{figure}
\centering
\includegraphics[width=0.4\textwidth]{fig4_03.pdf}
\caption{磁性氧化物中的超交换。箭头表示四个电子的自旋及其在过渡金属（M）与氧（O）
原子上的分布。设 M 有一个未成对电子因而具有磁性。若金属原子上的磁矩反铁磁耦合
（a、b、c），基态为（a），它可与（b）、（c）这样的激发组态混合：磁性电子因而在
M--O--M 单元上离域化，动能降低。若金属原子上的磁矩铁磁耦合（d、e、f），基态（d）
无法与（c）、（f）这样的激发组态混合——不相容原理禁止这些组态。铁磁组态因此付出
更多能量。}
\end{figure}

由于超交换同时涉及氧轨道与金属原子，它是二级过程，由二阶微扰论给出。二阶微扰论的
一般结论是：所涉能量近似等于跃迁矩阵元的平方除以制造激发态的能量代价。这里跃迁
矩阵元由一个称为跳跃积分（hopping integral）$t$ 的参数控制——在简单紧束缚近似中
它正比于导带的能量宽度（即带宽）。制造激发态的能量代价由库仑能 U 给出。于是
$\mathsf{J}\sim-t^{2}/U$。（四阶时可以有 $\Delta E\propto-t^{4}(\mathbf{S}_1\cdot\mathbf{S}_2)^{2}/U^{3}$
形式的相互作用，称为双二次交换（biquadratic exchange）。）

交换积分由两部分组成。第一是势交换项，代表电子排斥，更倾向于铁磁基态，但在离子
相距较远时很小。第二是动交换项，在此占主导，即上述效应。它依赖于轨道交叠程度，
因此超交换强烈依赖 M--O--M 键角。图只为一种 d 轨道而画；能与氧轨道交叠的其他 d
轨道的效应可能也需计入。

在某些情形下，超交换实际上可以是铁磁的。设想这样一个情形：通过一个氧离子，一个
磁性离子上的已占据 $e_g$ 轨道与另一磁性离子上的未占据 $e_g$ 轨道相耦合。$e_g$
电子跳到未占据轨道上有能量好处——只要到达时其自旋因洪特规则耦合而与 $t_{2g}$
电子的自旋排齐。此时超交换可以是铁磁的，但这是较弱的相互作用，也比常见的反铁磁
超交换少见。

\subsection{金属中的间接交换}

金属中磁性离子之间的交换相互作用可以由传导电子传递。局域磁矩使传导电子自旋极化，
这一极化又与距离 $r$ 之外的另一个局域磁矩耦合。这种交换是间接的——不涉及磁矩间的
直接耦合——称为 RKKY 相互作用（也称巡游交换，itinerant exchange）。RKKY 之名取自
该效应发现者 Ruderman、Kittel、Kasuya、Yosida 四人姓氏的首字母。耦合取与 $r$ 相关
的交换相互作用 $\mathsf{J}_{\mathrm{RKKY}}(r)$ 的形式：
\begin{equation*}
\mathsf{J}_{\mathrm{RKKY}}(r) \propto \frac{\cos(2k_{\mathrm{F}}r)}{r^{3}}.\tag{4.12}
\end{equation*}
（大 $r$ 时，设半径 $k_{\mathrm{F}}$ 的球形费米面。）这一相互作用是长程的，且随磁矩
间距振荡：依间距不同，可以铁磁也可以反铁磁。由于费米面十分锐利，耦合以 $\pi/k_{\mathrm{F}}$
为波长振荡。RKKY 相互作用将在第七章详细讨论。

\subsection{双交换}

在某些氧化物中可以出现铁磁交换相互作用，其原因是磁性离子可表现混合价态
（mixed valency）——即以不止一种氧化态存在。例子包括含 Mn 离子的化合物：Mn 可
处于氧化态 3 或 4，即 $\mathrm{Mn^{3+}}$ 或 $\mathrm{Mn^{4+}}$。其一即
$\mathrm{La}_{1-x}\mathrm{Sr}_x\mathrm{MnO}_3$（$0\leq x\leq1$），采取钙钛矿结构。
Sr 是二价的（以 $\mathrm{Sr^{2+}}$ 存在），La 是三价的（以 $\mathrm{La^{3+}}$ 存在）。
这意味着有 $x$ 分数的 Mn 离子是 $\mathrm{Mn^{4+}}$，$1-x$ 是 $\mathrm{Mn^{3+}}$。
系列两端（$x=0$ 与 $x=1$）都是反铁磁绝缘体——对磁性经氧由超交换传递的氧化物材料
而言这正合预期。$\mathrm{LaMnO_3}$ 只含 $\mathrm{Mn^{3+}}$，而 $\mathrm{Mn^{3+}}$
是扬--特勒离子；$\mathrm{LaMnO_3}$ 具有 A 型反铁磁有序（见 5.2 节）。然而当
$\mathrm{LaMnO_3}$ 掺 Sr 至 $x=0.175$ 时，扬--特勒畸变消失，系统变为铁磁，居里
温度在室温附近；低于该温度材料变为金属。

铁磁排列源于双交换（double exchange）机制，可借图 4.4 理解。$\mathrm{Mn^{3+}}$ 离子
上的 $e_g$ 电子只有当邻位存在同自旋空位时才能跳过去（因为跳跃不改变跳跃电子的
自旋）。若邻居是 $e_g$ 壳层无电子的 $\mathrm{Mn^{4+}}$，这本不成问题。但 $e_g$
电子与 $t_{2g}$ 能级三个电子之间存在强单心（洪特第一规则）交换相互作用，要把它们
全部排齐。因此，$e_g$ 电子跳到其 $t_{2g}$ 自旋将与自己反平行的邻居离子上，在能量上
是不利的（图 4.4(b)）。要维持给体与受体离子上的高自旋排列，相邻离子必须铁磁排列。
由于跳跃能力带来动能节省，允许图 4.4(a) 所示的跳跃过程降低总能量：系统于是铁磁
排列以省能量。而且铁磁排列进而允许 $e_g$ 电子在晶体中穿行，材料变为金属。双交换
铁磁体的导电性问题将在 8.9.5 节进一步讨论。双交换本质上是扩展体系中的铁磁超交换。
\fn{2}{铁磁超交换通常用于两个孤立离子：铁磁排列省下的动能对应跳入激发态。双交换用于扩展体系：省下的动能对应电子带宽的增益。}

\begin{figure}
\centering
\includegraphics[width=0.55\textwidth]{fig4_04.pdf}
\caption{双交换机制：参与电子转移的 $\mathrm{Mn^{3+}}$（$3\mathrm{d}^{4}$）与
$\mathrm{Mn^{4+}}$（$3\mathrm{d}^{3}$）离子之间产生铁磁耦合。单心交换相互作用
有利于跳跃，当（a）相邻离子铁磁排列时；（b）相邻离子反铁磁排列时则否。}
\end{figure}

双交换也见于磁铁矿（$\mathrm{Fe_3O_4}$）：八面体位上有等量的 $\mathrm{Fe^{2+}}$
（$3\mathrm{d}^{6}$）与 $\mathrm{Fe^{3+}}$（$3\mathrm{d}^{5}$）离子，四面体位上另有同等数目的
$\mathrm{Fe^{3+}}$ 离子。双交换相互作用使八面体位上的 $\mathrm{Fe^{2+}}$ 与
$\mathrm{Fe^{3+}}$ 离子铁磁排列；四面体位上的 $\mathrm{Fe^{3+}}$ 不参与这一相互
作用，与八面体位上的 $\mathrm{Fe^{3+}}$ 经反铁磁超交换耦合。于是两组
$\mathrm{Fe^{3+}}$ 的磁矩相消，剩下仅由 $\mathrm{Fe^{2+}}$ 贡献的净磁矩。实测的
每化学式单元磁矩非常接近仅由 $\mathrm{Fe^{2+}}$ 给出的预期值 $4\mu_{\mathrm{B}}$。

\subsection{各向异性交换相互作用}

自旋--轨道相互作用也可能以类似超交换中氧原子的方式扮演角色：此时激发态与氧无关，
而由磁性离子之一的自旋--轨道相互作用产生。于是一个离子的激发态与另一个离子的
基态之间存在交换相互作用。这称为各向异性交换相互作用（anisotropic exchange
interaction），又称 Dzyaloshinsky--Moriya 相互作用。作用于两个自旋 $\mathbf{S}_1$、
$\mathbf{S}_2$ 之间时，它在哈密顿量中给出 $\hat{\mathcal{H}}_{\mathrm{DM}}$：
\begin{equation*}
\hat{\mathcal{H}}_{\mathrm{DM}} = \mathbf{D}\cdot\mathbf{S}_1\times\mathbf{S}_2.\tag{4.13}
\end{equation*}
当晶体场对两磁离子连线中点具有反演对称时矢量 D 为零。但一般 D 可以不为零，此时
依对称性它平行或垂直于连接两自旋的连线。该相互作用的形式迫使 $S_1$ 与 $S_2$ 在
垂直于 D 的平面内接近直角，且取向使能量为负。其效果常常是使自旋倾摆（cant，即
略微转动一个小角）。它常见于反铁磁体，产生垂直于反铁磁自旋轴的小铁磁矩分量，
称为弱铁磁性（weak ferromagnetism）。例如 $\alpha$-$\mathrm{Fe_2O_3}$、
$\mathrm{MnCO_3}$ 与 $\mathrm{CoCO_3}$ 中即有此效应。

\begin{figure}
\centering
\includegraphics[width=0.5\textwidth]{fig4_05.pdf}
\caption{磁矩用相邻格点 $i$、$j$ 上的约化矩 $m_i$、$m_j$ 表示，两点相距
$r_{ij}$，磁矩间夹角为 $\phi_{ij}$。约化矩 $m_i$、$m_j$ 按定义为单位矢量。}
\end{figure}

\subsection{连续介质近似}

本节回到式 4.7 的海森伯模型。对本书后文有用的是：在忽略晶格分立性的连续介质近似
下，为这一相互作用找出一个表达式。

先设 $\mathsf{J}_{ij}$ 在 $i$、$j$ 为最近邻时等于常数 $\mathsf{J}$，否则为零。写成
\begin{equation*}
\hat{\mathcal{H}} = -\sum_{\langle ij\rangle}\mathsf{J}\,\mathbf{S}_i\cdot\mathbf{S}_j,\tag{4.14}
\end{equation*}
求和号下的 $\langle ij\rangle$ 表示只对最近邻求和。考虑经典自旋，设最近邻自旋间
夹角为 $\phi_{ij}$ 且极小，即对所有 $i$、$j$ 有 $\phi_{ij}\ll1$。我们实质上是假定系统
表现铁磁性（见下一章）但自旋并未完全排齐。

在这些假设下，系统能量可写为
\begin{equation*}
E = -\mathsf{J}S^{2}\sum_{\langle ij\rangle}\cos\phi_{ij} = \text{常数} + \frac{\mathsf{J}S^{2}}{2}\sum_{\langle ij\rangle}\phi_{ij}^{2},\tag{4.15}
\end{equation*}
最后一步用了 $\phi_{ij}\ll1$ 时的 $\cos\phi_{ij}\approx1-\phi_{ij}^{2}/2$。下面
忽略常数项（它只是完全排齐态的能量）。定义约化矩 $\mathbf{m}=\mathbf{M}/M_{\mathrm{s}}$，
$M$ 为磁化强度，$M_{\mathrm{s}}$ 为饱和磁化强度。单位矢量 $\mathbf{m}$ 跟随自旋方向，
$m_x$、$m_y$、$m_z$ 可视作格点 $\mathbf{r}_{ij}$ 处自旋的方向余弦。用图 4.5 的记号，
可写
\begin{equation*}
|\phi_{ij}| \approx |\mathbf{m}_i-\mathbf{m}_j| \approx |(\mathbf{r}_{ij}\cdot\nabla)\mathbf{m}|,\tag{4.16}
\end{equation*}
于是能量可写为
\begin{equation*}
E = \mathsf{J}S^{2}\sum_{\langle ij\rangle}[(\mathbf{r}_{ij}\cdot\nabla)\mathbf{m}]^{2}.\tag{4.17}
\end{equation*}
连续介质极限下忽略晶格分立性，写成
\begin{equation*}
E = A\int_{V}[(\nabla m_x)^{2}+(\nabla m_y)^{2}+(\nabla m_z)^{2}]\,\mathrm{d}^{3}r,\tag{4.18}
\end{equation*}
其中
\begin{equation*}
A = 2\mathsf{J}S^{2}z/a,\tag{4.19}
\end{equation*}
$a$ 是最近邻距离，$z$ 是单胞内的格点数（简单立方 $z=1$，体心立方 bcc $z=2$，面心立
方 fcc $z=4$）。

这个结果也可以从另一视角得到。若断言交换来自非均匀的磁化分布，那么当非均匀性
较为平缓时，可以仅凭对称性导出结果：表达式必须在自旋旋转下、以及在磁化分量变号下
不变。因此寻找与晶体对称性相容的、M 的导数的最低偶数阶表达式。由于正比于
$\partial M_{\alpha}/\partial x_{\beta}$ 的项在 M 反向下变号，只剩磁化梯度的二次项。
最一般的表达式是
\begin{equation*}
E = \sum_{\alpha\beta\gamma}C_{\alpha\beta}\frac{\partial M_{\gamma}}{\partial x_{\alpha}}\frac{\partial M_{\gamma}}{\partial x_{\beta}},\tag{4.20}
\end{equation*}
$C_{\alpha\beta}$ 是具有晶体对称性的张量。立方晶体中它约化为
\begin{equation*}
E = C\sum_{\alpha\gamma}\frac{\partial M_{\gamma}}{\partial x_{\alpha}}\frac{\partial M_{\gamma}}{\partial x_{\alpha}} = C\left[(\nabla M_x)^{2}+(\nabla M_y)^{2}+(\nabla M_z)^{2}\right],\tag{4.21}
\end{equation*}
与上面式 4.18 的结果等价。式 4.21 有时写作 $C|\nabla\mathbf{M}|^{2}$，但要注意对
矢量取梯度须小心。

\section*{延伸阅读}

\begin{itemize}[itemsep=0.2em]
\item 关于交换相互作用的更多信息见 C.~Herring，载于《Magnetism》，G.~Rado 与
H.~Suhl 编，vol.~2B, p.1，Academic Press, New York 1966。
\item 磁性体系中相互作用的理论论述见 K.~Yosida，《Theory of magnetism》，
Springer 1996。
\item 真实体系中各种交换相互作用的综述见 P.~A.~Cox，《Transition metal oxides》，
OUP 1995。
\item 关于交换与交换相互作用非常透彻而有益的参考书是 D.~C.~Mattis，《The theory
of magnetism I》，Springer 1981。
\end{itemize}

\section*{习题}

\exer{4.1} 证明相距 $\mathbf{r}$ 的两个磁偶极子 $\boldsymbol{\mu}_1$ 与
$\boldsymbol{\mu}_2$ 的偶极能量为
\begin{equation*}
E = \frac{\mu_0}{4\pi r^{3}}\left[\boldsymbol{\mu}_1\cdot\boldsymbol{\mu}_2 - \frac{3}{r^{2}}(\boldsymbol{\mu}_1\cdot\mathbf{r})(\boldsymbol{\mu}_2\cdot\mathbf{r})\right].\tag{4.22}
\end{equation*}

\exer{4.2} 计算距质子 1 \AA{} 与 10 \AA{} 处磁场的大小，方向分别（a）平行与
（b）垂直于质子自旋方向。

\exer{4.3} 估计金属 Fe 中相邻两个 Fe 原子的交换耦合与偶极耦合之比。（Fe 的交换
常数可粗略地取 $k_{\mathrm{B}}T_{\mathrm{C}}$，$T_{\mathrm{C}}$ 为居里温度；Fe 的
$T_{\mathrm{C}}=1043$ K。）

\exer{4.4} 用（非弹性中子散射测得的）电子带宽 0.05 eV 粗略估计由超交换耦合的磁性
氧化物中交换常数的大小。取库仑能 $\sim$1 eV。进而估计反铁磁有序温度。

\exer{4.5} 考虑两个相互作用的自旋 $\frac{1}{2}$ 电子。好量子数为 $S=0$ 与 $1$，
故有一个三重态和一个单态，二者被能隙 $\Delta$ 分开。定义 $\Delta$ 的符号使
$\Delta>0$ 时单态（$S=0$）在下、$\Delta<0$ 时三重态在下。两种情形示于图 4.6(a)、
(b)。证明该模型的磁化率为
\begin{equation*}
\chi = \frac{2Ng\mu_{\mathrm{B}}^{2}}{k_{\mathrm{B}}T(3+\mathrm{e}^{\Delta/k_{\mathrm{B}}T})},\tag{4.23}
\end{equation*}
即 Bleaney--Bowers 方程，绘于图 4.6(c)。

\begin{figure}
\centering
\includegraphics[width=0.42\textwidth]{fig4_06.pdf}
\caption{由 $\Delta$ 耦合的两个自旋产生单态与三重态：（a）$\Delta>0$ 与（b）
$\Delta<0$ 时的能级。（c）由 Bleaney--Bowers 方程给出的磁化率，分别对 $\Delta>0$
（下曲线）、$\Delta=0$（中曲线）与 $\Delta<0$（上曲线）绘制。}
\end{figure}

\exer{4.6} 两个自旋为 $\mathbf{I}_1$、$\mathbf{I}_2$ 的核相距 $\mathbf{r}$，由偶极--
偶极相互作用耦合，哈密顿量为
\begin{equation*}
\hat{\mathcal{H}} = \frac{\mu_0}{4\pi r^{3}}\left[\hat{\boldsymbol{\mu}}_1\cdot\hat{\boldsymbol{\mu}}_2 - \frac{3}{r^{2}}(\hat{\boldsymbol{\mu}}_1\cdot\mathbf{r})(\hat{\boldsymbol{\mu}}_2\cdot\mathbf{r})\right],\tag{4.24}
\end{equation*}
其中 $\boldsymbol{\mu}_i=g_I\mu_{\mathrm{N}}\mathbf{I}_i$，$i=1,2$。证明哈密顿量可表示为
\begin{equation*}
\hat{\mathcal{H}} = \frac{\mu_0 g_I^{2}\mu_{\mathrm{N}}^{2}}{4\pi r^{3}}\left[A+B+C+D+E+F\right],\tag{4.25}
\end{equation*}
其中
\begin{equation*}
\begin{aligned}
A &= (1-3\cos^{2}\theta)I_{1z}I_{2z}\\
B &= -\frac{1}{4}(1-3\cos^{2}\theta)(I_{1}^{+}I_{2}^{-}+I_{1}^{-}I_{2}^{+})\\
C &= -\frac{3}{2}\sin\theta\cos\theta\,e^{-i\phi}(I_{1z}I_{2}^{+}+I_{1}^{+}I_{2z})\\
D &= -\frac{3}{2}\sin\theta\cos\theta\,e^{i\phi}(I_{1z}I_{2}^{-}+I_{1}^{-}I_{2z})\\
E &= -\frac{3}{4}\sin\theta\,e^{-2i\phi}I_{1}^{+}I_{2}^{+}\\
F &= -\frac{3}{4}\sin\theta\,e^{2i\phi}I_{1}^{-}I_{2}^{-},
\end{aligned}\tag{4.26}
\end{equation*}
角 $\theta$ 与 $\phi$ 是矢量 $\mathbf{r}$ 的球极角，升降算符为
$I_j^{\pm}=I_{jx}\pm\mathrm{i}I_{jy}$。对三重态（即 $\mathbf{I}=\mathbf{I}_1+\mathbf{I}_2$
大小 $I=1$）的两个质子（$I_1=I_2=\frac{1}{2}$），证明上列哈密顿量中的各项诱导如下
跃迁：$A$ 与 $B$ 为 $\Delta M=0$，$C$ 为 $\Delta M=+1$，$D$ 为 $\Delta M=-1$，
$E$ 为 $\Delta M=+2$，$F$ 为 $\Delta M=-2$。

\exer{4.7} 两个原子键合形成双原子分子。分子中有两个电子，坐标为 $\mathbf{r}_1$、
$\mathbf{r}_2$；两个核分别固定在 $\mathbf{R}_1$、$\mathbf{R}_2$，电荷各为 $Ze$。
忽略核间排斥，哈密顿量可写为 $\hat{\mathcal{H}}_0+\hat{\mathcal{H}}'$，其中
\begin{equation*}
\hat{\mathcal{H}}_0 = -\frac{\hbar^{2}}{2m}(\nabla_1^{2}+\nabla_2^{2}) - \frac{Ze^{2}}{4\pi\epsilon_0|\mathbf{r}_1-\mathbf{R}_1|} - \frac{Ze^{2}}{4\pi\epsilon_0|\mathbf{r}_2-\mathbf{R}_2|} = \hat{\mathcal{H}}_1+\hat{\mathcal{H}}_2,\tag{4.27}
\end{equation*}
是单粒子哈密顿量 $\hat{\mathcal{H}}_1$ 与 $\hat{\mathcal{H}}_2$（每个电子绕自己的
核）之和，而
\begin{equation*}
\hat{\mathcal{H}}' = \frac{e^{2}}{4\pi\epsilon_0|\mathbf{r}_1-\mathbf{r}_2|} - \frac{Ze^{2}}{4\pi\epsilon_0|\mathbf{r}_1-\mathbf{R}_2|} - \frac{Ze^{2}}{4\pi\epsilon_0|\mathbf{r}_2-\mathbf{R}_1|},\tag{4.28}
\end{equation*}
包含电子排斥能与每个电子受对方核吸引的能量项。证明分子能量为
\begin{equation*}
E = E_1 + E_2 + \frac{K\pm\mathsf{J}}{1\pm S}\tag{4.29}
\end{equation*}
$E_1$、$E_2$ 分别是 $\hat{\mathcal{H}}_1$、$\hat{\mathcal{H}}_2$ 的本征值，
\begin{equation*}
K = \int\psi_a^{*}(\mathbf{r}_1)\psi_b^{*}(\mathbf{r}_2)\hat{\mathcal{H}}'\psi_a(\mathbf{r}_1)\psi_b(\mathbf{r}_2)\,\mathrm{d}\mathbf{r}_1\mathrm{d}\mathbf{r}_2 = \int|\psi_a(\mathbf{r}_1)|^{2}\hat{\mathcal{H}}'|\psi_b(\mathbf{r}_2)|^{2}\,\mathrm{d}\mathbf{r}_1\mathrm{d}\mathbf{r}_2\tag{4.30}
\end{equation*}
称为库仑积分，
\begin{equation*}
\mathsf{J} = \int\psi_a^{*}(\mathbf{r}_1)\psi_b^{*}(\mathbf{r}_2)\hat{\mathcal{H}}'\psi_a(\mathbf{r}_2)\psi_b(\mathbf{r}_1)\,\mathrm{d}\mathbf{r}_1\mathrm{d}\mathbf{r}_2\tag{4.31}
\end{equation*}
称为交换积分，
\begin{equation*}
S = \int\psi_a^{*}(\mathbf{r}_1)\psi_b^{*}(\mathbf{r}_2)\psi_a(\mathbf{r}_2)\psi_b(\mathbf{r}_1)\,\mathrm{d}\mathbf{r}_1\mathrm{d}\mathbf{r}_2\tag{4.32}
\end{equation*}
称为交叠积分，且假定本征函数为
\begin{equation*}
\Psi_{\pm} = \frac{1}{\sqrt{2}}\left[\psi_a(\mathbf{r}_1)\psi_b(\mathbf{r}_2)\pm\psi_a(\mathbf{r}_2)\psi_b(\mathbf{r}_1)\right].\tag{4.33}
\end{equation*}

\exer{4.8} 三个 $S=1$ 原子置于等边三角形的三个顶点，可用哈密顿量
\begin{equation*}
\hat{\mathcal{H}} = -2\mathsf{J}(\hat{\mathbf{S}}_1\cdot\hat{\mathbf{S}}_2+\hat{\mathbf{S}}_2\cdot\hat{\mathbf{S}}_3+\hat{\mathbf{S}}_3\cdot\hat{\mathbf{S}}_1).\tag{4.34}
\end{equation*}
描述。证明系统的能量本征值为 0、$2\mathsf{J}$ 与 $6\mathsf{J}$。
